\documentclass[conference]{IEEEtran}
\usepackage[utf8]{inputenc}
\usepackage[T1]{fontenc}
\usepackage{cite}
\usepackage{amsmath,amssymb}
\usepackage{graphicx}
\usepackage{booktabs}
\usepackage{url}

\title{Verifier-Guided Prompting for Intent-to-Configuration Translation: A Preregistered NetOps Study}

\author{
\IEEEauthorblockN{Autonomous AI Researcher}
\IEEEauthorblockA{CNSM 2026 Agentic AI Researcher Track}
}

\begin{document}

\maketitle

\begin{abstract}
We preregistered and executed a paired confirmatory trial to test whether constrained prompting plus a deterministic verifier-guided generate\ensuremath{\rightarrow}verify\ensuremath{\rightarrow}repair loop (one allowed repair) increases deterministic-verifier semantic-correctness when translating operator natural-language intents into low-level device configurations. Design: N=150 synthetic intents sampled from a deterministic generator, shared\_initial\_candidate generation semantics, model = gpt-5-mini, deterministic validator = deterministic\_netops\_validator\_v5, one initial generation per intent shared across baseline and guarded conditions, guarded pipeline permitted <=1 repair call. Results (artifact-grounded): baseline = 149/150 (99.333\%); guarded = 150/150 (100.0\%); paired absolute difference = 0.006667 (\ensuremath{\approx}0.67 percentage points); exact McNemar two-sided p = 1.0; 95\% bootstrap percentile CI = [0.00, 0.02] (10,000 resamples, seed = 1729). Provider accounting: completed episodes = 300; model\_calls\_used = 151; repair-stage invocations = 1. Interpretation: on this preregistered synthetic corpus and under shared\_initial\_candidate semantics the guarded workflow produced a numerically negligible and statistically non-significant improvement over a near-ceiling baseline. We emphasize the methodological lesson that preregistered NetOps trials require baseline-calibration pilots to avoid ceiling effects that invalidate preregistered power assumptions. All execution and analysis artifacts are archived in the supplied execution bundle referenced by the execution manifest.
\end{abstract}

\section{Introduction}
Natural-language intent interfaces aim to reduce operator burden by letting humans express desired network changes as free text that is automatically translated into executable device configurations. Prior work documents that single-pass LLM outputs often contain syntactic errors, omissions, or semantic violations that can yield unsafe or unavailable states if applied directly [1][2][3]. A commonly proposed defense is a generate\ensuremath{\rightarrow}verify\ensuremath{\rightarrow}repair architecture in which an LLM produces a candidate, a deterministic verifier checks intent-level invariants or formalized constraints, and an automated repair (or human gate) corrects violations before execution [4][5]. Security studies of tool-augmented LLM agents further motivate deterministic gating because unchecked textual claims can become harmful actions in multi-tool workflows [6].

Research question. After an autonomous literature review and preregistration we posed a confined confirmatory question: does constrained prompting (structured templates with explicit semantic slots) combined with deterministic verifier-guided iterative feedback materially increase verifier-level semantic correctness of NL\ensuremath{\rightarrow}configuration translation versus an unconstrained baseline under a paired design? The trial (study\_cand\_2026\_b2\_v1) was preregistered and frozen before any model outputs were observed.

Contributions. (1) A fully preregistered, artifactized paired confirmatory trial measuring deterministic-verifier pass/fail on a programmatically generated synthetic corpus (N=150 paired intents). (2) Artifact-backed outcomes showing that, under shared\_initial\_candidate semantics and a one-repair budget, the guarded pipeline raised verifier pass rate from 149/150 to 150/150 (+0.67 pp), a numerically negligible and statistically non-significant improvement (exact McNemar p = 1.0; 95\% bootstrap CI = [0.00, 0.02]). (3) A methodological recommendation that preregistered NetOps trials include baseline-calibration pilots and explicitly specify generation semantics (shared\_initial\_candidate vs independent sampling) because semantics materially change the estimand and interpretation.

\section{Related Work}
Related literature falls into three strands; each strand provides partial motivation for the guarded generate\ensuremath{\rightarrow}verify\ensuremath{\rightarrow}repair architecture we evaluate and clarifies where prior claims and limitations remain.

Intent-to-configuration translation. A growing body of work explores converting operator natural-language intents into formal specifications or concrete device configurations. Empirical evaluations and surveys report that single-pass translation often attains only modest semantic-correctness on curated benchmarks, motivating post-generation checks and constrained prompting to improve fidelity [1][2]. Representative studies highlight varied task formulations (intent grammars, template-based translation, and intent-slot extraction) and report that success rates depend strongly on task granularity, the expressiveness of the target formalism, and prompt design choices [1]. Systematic reviews and targeted evaluations underscore that many NL\ensuremath{\rightarrow}formal pipelines still produce semantic errors that formal verifiers can detect but that the end-to-end reliability required for unattended NetOps remains an open problem [2]. These findings justify using deterministic verifiers as the primary confirmatory endpoint and motivate task-bank designs that capture common operational idioms (reachability, ACL-like constraints, VLAN/MTU sequences, and uplink switchover patterns) rather than only isolated syntactic transforms.

Generate\ensuremath{\rightarrow}verify\ensuremath{\rightarrow}repair architectures and verifier technology. Several recent method papers and systems integrate deterministic or policy-guided verifiers with generation to improve correctness for infrastructure-as-code and configuration-repair tasks [4][7]. Approaches range from using verifier feedback for fine-tuning to runtime repair loops that iteratively correct a candidate until invariants hold. Reported improvements are heterogeneous: gains are larger when baseline model competence is modest, when larger repair budgets or multiple independent samples are allowed, or when verifier feedback provides actionable, fine-grained corrections rather than coarse pass/fail signals [4][7]. Critically, the effectiveness of verifier-guided repair depends on verifier coverage and the fidelity of the mapping layer that translates natural-language intents to verifier inputs; prior work emphasizes rigorous unit-testing of mapping code and staged integration with formal tools [5].

A complementary thread documents mature formal verification tools (e.g., NetKAT variants, KATch, Hoyan, and other symbolic verifiers) that can act as deterministic oracles for reachability, isolation, and policy checks in controlled topologies [5][9]. These tools vary in supported abstractions, latency, and scale: some are optimized for design-time proofs while others (e.g., recent NetKAT provers and scalable overlay checkers) demonstrate faster runtime checks for smaller topologies. Prior studies therefore treat verifier selection and verifier-coverage testing as design decisions that directly affect measured semantic-correctness and external validity.

Security, prompt-injection, and claim-to-action risks. Security- and attack-oriented analyses show that LLM-driven multi-tool workflows can convert false textual claims into concrete tool calls and harmful actions unless guarded by auditable gates or hardened planning primitives [6][8]. Empirical prompt-injection studies demonstrate practical exploitation strategies and optimization-based attacks against LLM-as-judge setups; corresponding defenses (prompt hardening, verifier-aware gating, and call-auditing) reduce risk but impose trade-offs with automation throughput [6][8]. These results motivate conservative guarded architectures that retain verifier traces and normalized configurations for operator review, a design choice we make explicit when framing operational interpretation beyond raw pass/fail rates.

How this study connects and differs. The present trial operationalizes a narrowly specified guarded architecture informed by the above strands: we evaluate verifier-level pass/fail on a deterministic synthetic corpus (so verifier selection and mapping tests are central), we limit repair budget to a single repair call to match a constrained operational budget, and we adopt shared\_initial\_candidate semantics to isolate the marginal effect of repair applied to an identical first-pass candidate. These design choices intentionally expose how sampling semantics, verifier coverage, and repair budgets interact with baseline competence---issues emphasized across the cited literature---and explain why previously reported larger guarded-vs-baseline gains can co-exist with a near-ceiling outcome in tightly controlled, verifier-capable synthetic settings [4][5][7].

\section{Methodology}
Preregistration and confirmatory design. The study was preregistered as study\_cand\_2026\_b2\_v1 with a frozen analysis plan executed before any model outputs were observed. The primary estimand was the paired\_success\_rate\_difference\_guarded\_minus\_baseline evaluated with an exact two-sided McNemar test on deterministic-verifier pass/fail outcomes. The confirmatory sample specified N=150 distinct synthetic intents drawn deterministically from a preregistered generator, deterministic inclusion/exclusion rules (a rule-based ambiguity detector and verifier-coverage checks), a deterministic retry policy for provider timeouts (one retry), and a power target to detect an absolute +15 percentage-point improvement (\ensuremath{\approx}80\% power) under conservative discordant-pair assumptions. All elements of the preregistration (estimands, inclusion/exclusion rules, multiplicity plan, and analysis scripts) were frozen and archived prior to model calls.

Execution semantics: shared\_initial\_candidate, call budgets, and adapter enforcement. A central design decision was the execution-semantic enforcement of shared\_initial\_candidate: for each intent the execution adapter produced exactly one initial model generation and that identical sampled candidate was used in the paired evaluation of both baseline and guarded conditions. The guarded pipeline was permitted at most one additional repair call only if the deterministic verifier failed that initial candidate. Operationally, this means the baseline rate measures the validity of the single initial sampled candidate; the guarded vs baseline contrast therefore isolates the marginal effect of applying at most one repair to that identical candidate rather than measuring performance averaged over independent first-pass draws. The adapter enforced these constraints at runtime and recorded stage-level provider accounting; archived provider audits report stage\_counts = \{shared\_initial\_generation: 150, repair: 1\} and model\_calls\_used = 151, consistent with the enforced one-initial plus at most-one-repair policy.

Model, sampling, and deterministic validator. Declared model: gpt-5-mini (provider record: openai\_responses). The execution metadata include an explicit record that model\_configuration.json declares declared\_model\_version = "gpt-5-mini" and temperature = null (unset). We emphasize the literal meaning: temperature = null in the recorded configuration is not equivalent to setting temperature = 0 (deterministic decoding). Because no numeric temperature value or fixed random-seed was enforced in the preregistration or execution, model outputs may be nondeterministic across independent API sessions. This nondeterminism was anticipated in the preregistration and is reported explicitly here because sampling semantics materially affect estimand interpretation. Deterministic verifier: deterministic\_netops\_validator\_v5 (validator\_version = 5.0) applied a full\_intent\_constraint\_validation rule and returned for each candidate a boolean pass/fail verdict plus normalized final-configuration text and per-episode diagnostics (parsed\_command\_count, required\_command\_count, observed\_touched\_field\_count, violation\_codes). These deterministic verdicts constituted the binary endpoint for the confirmatory test.

Task generator, corpus, prechecks, and mapping verification. Confirmatory tasks were produced deterministically by the preregistered generator (deterministic\_netops\_task\_generator\_v5) and were stratified across four intent families: reachability/isolation, ACL-like access changes, VLAN/MTU sequence changes, and uplink switchover patterns. Topologies were constrained to verifier-capable small networks (4--32 nodes) to ensure deterministic oracle coverage. Prior to batch execution we ran the preregistered mapping and verifier-coverage unit-tests: parser transformations that map LLM textual outputs to verifier inputs were unit-tested and passed. The preregistration specified deterministic exclusion rules (ambiguous intents routed to exploratory corpus; parse failures or unsupported verifier features excluded with deterministic reason codes); no confirmatory exclusions occurred in the archived run. Representative generator seeds, task\_payloads, and expected\_final\_states are archived to support independent calibration or pilot reuse.

Scoring rules, primary analysis, and confidence procedures. The primary outcome was deterministic verifier pass (binary). Scoring used the archival scorer identifier deterministic\_netops\_validator\_v5 and its full\_intent\_constraint\_validation rule; for each episode the scorer produced a normalized\_configuration string that was compared to the preregistered reference invariants for automated pass/fail determination. The confirmatory analysis used an exact two-sided McNemar test on the paired contingency (n\_11, n\_10, n\_01, n\_00) and computed a paired bootstrap percentile 95\% confidence interval with 10,000 resamples (bootstrap\_seed = 1729). The preregistered handling of incomplete pairs specified exclusion from the primary test (complete\_pair\_primary). Secondary contrasts (constrained vs baseline, guarded vs constrained) and multiplicity corrections were preregistered; under shared\_initial\_candidate enforcement these secondary contrasts were constrained by the adapter semantics and the available stage counts. All analysis code, contingency tables, and bootstrap parameters are archived and cited in the analysis artifacts.

Execution accounting, contamination controls, and reproducibility. Planned episodes were 300 (150 intents \ensuremath{\times} 2 conditions); completed\_episode\_count = 300 and complete\_pair\_count = 150 in the archived manifest. Provider-call accounting and contamination checks were recorded automatically: completed\_hosted\_model\_call\_event\_count = 151, failed\_hosted\_model\_call\_event\_count = 0, cache\_miss\_count = 151, and contamination\_summary.csv recorded zero flags. These runtime audits and per-call metadata (provider traces, shared\_initial\_response paths, and validator traces) were retained in the execution bundle to enable deterministic reconciliation of the paired counts and to support post hoc reproducibility checks. The preregistration also specified a contamination-detection checklist (session isolation, duplicate-response checks, mapping unit-tests) that the adapter enforced prior to batch execution; all checks passed and are documented in the artifact logs.

\section{Results}
Execution and provider audit. The run completed exactly as planned: completed\_episode\_count = 300 and complete\_pair\_count = 150. Provider-call accounting (from the archived execution manifest) shows model\_calls\_used = 151, completed\_hosted\_model\_call\_event\_count = 151, failed\_hosted\_model\_call\_event\_count = 0, cache\_miss\_count = 151, and stage\_counts = \{shared\_initial\_generation: 150, repair: 1\}. Contamination checks raised no flags and all unit-tests that gate mapping and verifier-input transformations passed prior to batch execution.

Primary contingency and per-condition rates. The paired contingency table (guarded \ensuremath{\times} baseline) observed in the archived analysis artifacts was: n\_11 = 149 (both-pass), n\_10 = 1 (guarded-only-pass), n\_01 = 0 (baseline-only-pass), and n\_00 = 0 (both-fail). These counts yield per-condition success totals baseline\_success\_count = 149/150 = 0.9933333333333333 and guarded\_success\_count = 150/150 = 1.0, and an absolute paired difference (guarded - baseline) = 0.00666666666666671 (\ensuremath{\approx}0.67 percentage points). The preregistered exact two-sided McNemar test on these paired counts produced p = 1.0. The paired bootstrap-percentile 95\% confidence interval computed with 10,000 resamples (bootstrap\_seed = 1729) is [0.00, 0.02]. All these numeric values are taken verbatim from the archived analysis files (analysis/paired\_contingency\_table.csv and analysis/condition\_summary.csv).

Repair-stage diagnostics and revision iterations. The guarded pipeline issued a repair-stage model call in exactly one episode (1/150 tasks), matching the provider audit (stage\_counts.repair = 1). The archived scoring traces indicate that this single repair invocation converted an initially verifier-failing candidate into a verifier-accepted normalized\_configuration; the validator logs include before/after normalized text, violation\_codes, and parsed\_command\_count fields for that episode. Across representative archived traces (tasks 000001--000006) the parser- and validator-level diagnostics show parsed\_command\_count values consistent with the required\_command\_count reported by the generator (examples in the archive show parsed\_command\_count values of 2, 4, and 6 for representative easy-to-hard cases). Consistent with the rarity of repairs, the median revision iterations per guarded episode = 0 (IQR = 0), and only the single discordant pair required a guarded repair under the preregistered one-repair budget.

Representative artifact-grounded examples. Representative archived task traces illustrate the variety of structural difficulty encoded by the task generator. For example, task-000001 (generator\_seed = 1) required a four-step shutdown\ensuremath{\rightarrow}change\ensuremath{\rightarrow}restore sequence; the shared-initial candidate and both condition responses matched the required four-command sequence and the validator reported parsed\_command\_count = 4, required\_command\_count = 4, observed\_touched\_field\_count = 3, and valid = true. Other representative traces (e.g., a two-command make-before-break uplink switchover and a six-command paired-MTU change) demonstrate that the corpus includes both relatively compact and more complex required sequences; these representative artifacts (responses and validator traces) are archived and cited in the execution manifest.

Sensitivity and missingness outcomes. The preregistered sensitivity analysis planned to treat missing guarded outputs as failures was not exercised because failed\_hosted\_model\_call\_event\_count = 0 and no unscorable episodes occurred. The missingness summary in the archive reports complete\_pairs = 150 and all related counts = 0 for failed or unscorable episodes.

Compact evidence tables and artifacts. Table 1 (paired contingency) and Table 2 (execution audit summary) in the manuscript reproduce the exact archived CSVs (analysis/paired\_contingency\_table.csv and analysis/condition\_summary.csv). The analysis artifacts also record bootstrap parameters (bootstrap\_resamples = 10000; bootstrap\_seed = 1729) and the paired-analysis log entry that marks the completion of the confirmatory analysis.

Interpretation grounded in archived artifacts. On this preregistered synthetic corpus and under shared\_initial\_candidate semantics the guarded workflow yielded a numerically small absolute improvement (+0.67 pp) over an already near-ceiling baseline. The exact McNemar p = 1.0 and the bootstrap CI that includes zero indicate no statistically supported confirmatory benefit under the preregistered estimand. The empirical facts driving this conclusion are: (a) an observed baseline success rate \ensuremath{\approx}99.33\% that left essentially no headroom for a large absolute increase, and (b) a single, rare repair invocation that produced the only discordant pair. These artifact-grounded observations motivate the methodological recommendation to perform baseline-calibration pilots prior to confirmatory sampling to avoid underpowered trials caused by ceiling effects. All numeric statements above and the diagnostic traces referenced are archived in the execution bundle and analysis artifacts cited in the execution manifest.

\section{Discussion}
We expand the discussion with concrete, artifact-grounded diagnostics, methodological implications, and operational interpretations that directly follow from the archived execution and analysis artifacts.

Ceiling effect, discordance structure, and inferential consequence. The contingency structure (n\_11=149, n\_10=1, n\_01=0) makes explicit why the exact McNemar test returns p=1.0: almost all paired episodes agree and only a single discordant pair favors the guarded condition. Under McNemar's paired-proportion logic the test's power derives from the number of discordant pairs; with only one discordant pair there is effectively no evidence for a systematic directional effect. The paired bootstrap CI ([0.00,0.02], bootstrap\_seed=1729, resamples=10000) quantifies the same intuition by showing that plausible values for the absolute improvement under resampling include zero and are bounded above by \ensuremath{\approx}2 percentage points. Both diagnostics are reported verbatim from analysis/paired\_contingency\_table.csv and analysis/condition\_summary.csv and therefore trace directly to archived artifacts rather than post hoc interpretation.

Mechanistic explanation via execution semantics. The shared\_initial\_candidate semantics materially reduce the pathway by which the guarded pipeline can improve outcomes: because the identical initial candidate is used to evaluate both baseline and guarded conditions, any potential advantage of the guarded pipeline depends entirely on the single allowed repair call converting that specific initial candidate into a verifier-pass. This is reflected in provider accounting: stage\_counts show shared\_initial\_generation=150 and repair=1, and model\_calls\_used=151. Practically, when baseline initial candidates are already highly likely to pass the verifier (here \ensuremath{\approx}99.33\%), the repair budget is rarely exercised and the guarded condition has almost no opportunity to improve aggregate pass rates. This mechanism explains the observed near-ceiling result and illustrates how sampling semantics and repair budgets interact with baseline competence to determine observable effect sizes.

Design consequences for preregistration and power calibration. The preregistration targeted detection of a 15 percentage-point absolute improvement with N=150. That calculation assumed substantially lower baseline competence and therefore a larger expected discordant mass. Our artifact-grounded outcome demonstrates the risk of committing to a confirmatory N without an empirical baseline check: baseline calibration reduces the chance of an underpowered or effectively vacuous confirmatory test caused by ceiling effects. Concretely, a practical pilot procedure (consistent with our recommendations and implementable from the archived deterministic task generator) is: sample 20--40 tasks from the planned generator with the intended prompt styles; measure baseline verifier-pass rate on those pilot tasks; if baseline >\textasciitilde{}85--90\% either increase task difficulty or enlarge N and/or adopt independent-sampling semantics or a larger repair budget. These procedural steps require no new experimental artifacts and can be implemented with the provided generator seeds and unit-tested mapping code archived in the execution bundle.

Operational affordances beyond the binary endpoint. Even when aggregate verifier-pass improvements are negligible, the guarded pipeline provides operationally valuable artifacts that the binary pass/fail metric does not capture. The deterministic validator returns structured diagnostics (normalized\_configuration strings, parsed\_command\_count, required\_command\_count, violation\_codes) for each episode; these are visible in validator traces for representative tasks (e.g., task-000001..000006). In practice these strings enable: (1) auditable diffs for operator review, (2) concise failure taxonomies to prioritize human triage, and (3) deterministic gating rules implementable as fail-stop or auto-remediation policies. The run-level audit (completed\_hosted\_model\_call\_event\_count=151; failed\_hosted\_model\_call\_event\_count=0) proves these artifacts were produced consistently across the sample and are therefore operationally available in the archived bundle for downstream tooling.

Trade-offs among sampling semantics, repair budgets, and estimands. Our evidence clarifies three non-identical estimands that practitioners must choose between and that should be explicit in preregistrations: (A) shared\_initial\_candidate marginal-repair effect (our estimand), (B) independent-sampling expected-performance difference (draw independent initial samples per condition), and (C) repair-robustness with expanded repair budgets (allow multiple iterative repairs). Estimand A is conservative with respect to repair effectiveness because it fixes the initial draw; estimand B typically yields larger observable differences when model output variance is high; estimand C quantifies diminishing returns from additional repair attempts. The archived artifacts make it possible to re-run variants B or C (subject to provider budget and adapter constraints) starting from the same deterministic task generator seeds and the stored model\_configuration metadata; this preserves reproducibility while acknowledging that different, equally defensible operational choices lead to different scientific answers.

Reproducibility notes tied to concrete artifact fields. Key numeric analysis parameters are archived and should be used by any reproducer: bootstrap\_resamples=10000 and bootstrap\_seed=1729 (analysis/analysis\_log.jsonl), model\_calls\_used=151 and stage\_counts (analysis/deterministic\_reconciliation.json), and declared\_model\_version="gpt-5-mini" with model\_configuration.json recording temperature=null (unset). Because temperature was unset and no fixed random seed was enforced, repeated independent provider sessions can produce nondeterministic outputs; reproductions seeking exact token-level matches must therefore control sampling parameters or accept equivalent statistical replications rather than byte-identical replays. All of these facts are documented in the execution manifest and model\_configuration.json and are cited here to avoid ambiguity about what archive-grounded reproduction entails.

Threats to validity revisited with artifact grounding. Internal validity: the adapter's session-isolation and provider-call audit (no cross-condition cache hits observed) reduce concerns about prompt-history leakage; these checks are recorded in deterministic\_reconciliation.json. Construct validity: deterministic validator pass/fail reflects the predeclared invariants from the generator's task\_payload; mapping-code unit-tests passed pre-execution but any gaps in verifier coverage remain a construct threat---this is why we limited claims to the verifier-capable invariants and reported representative validator violation fields when present. External validity: the single-model constraint (gpt-5-mini) and small-topology synthetic corpus restrict generalization; these are enumerated in the preregistration and in the limitations section and are supported by the model\_configuration and task-manifest artifacts. Conclusion validity: the mismatch between preregistered power assumptions and the observed baseline is a concrete illustration of how inferential conclusions depend on realized effect sizes and discordant-pair counts, both archived in the paired\_contingency\_table.csv.

Practical recommendations for practitioners and experimenters. (1) Always run a short baseline-calibration pilot using the final prompt templates and generator to estimate baseline competence and select N, difficulty, and repair budgets accordingly. (2) Make generation semantics explicit in preregistration: shared\_initial\_candidate vs independent sampling lead to different operational decisions and must match intended deployment semantics. (3) When verifier coverage is partial, report both binary pass/fail and the structured validator diagnostics to allow operators to triage and to enable metric-based comparisons across verifier designs. (4) Consider multi-arm designs that vary repair budgets and sampling semantics to estimate diminishing returns and to align statistical power with expected discordance rates. The provided execution bundle contains deterministic generator seeds, unit-tested mapping code, and validator traces sufficient to implement each recommendation without inventing new artifacts.

Summary. The expanded discussion ties our artifact-grounded numerical outcomes to concrete methodological mechanisms (sampling semantics and repair budget), diagnostic artifacts (validator traces and audit counts), reproducibility parameters (bootstrap seed, model metadata), and explicit experimental recommendations. These points are strictly supported by the archived execution and analysis artifacts and do not alter empirical claims reported elsewhere in the manuscript.

\section{Limitations}
\begin{itemize}
\item Single-model constraint: experiments used only gpt-5-mini; results do not generalize across models or sizes.
\item Synthetic corpus and scale: tasks were programmatically generated and limited to verifier-capable toy-to-small topologies (4--32 nodes); external validity to production WANs and heterogeneous vendor CLIs is untested.
\item Shared\_initial\_candidate semantics and execution contract limited repair budget to at most one guarded repair call per intent; different generation semantics or multiple-repair budgets may yield different outcomes.
\item Ceiling effect and preregistration mismatch: baseline correctness was far higher than the pilot assumptions used for power calculations (planned detectable effect = 15 percentage points), reducing the trial's ability to demonstrate benefit and illustrating sensitivity to baseline calibration.
\item Verifier coverage and specification translation: the study measures deterministic validator pass/fail on the preregistered invariants but does not claim safety guarantees beyond the deterministic validator's modeled properties.
\item Security and adversarial robustness: the experiment did not evaluate adversarial prompt-injection or adversary-driven intents; separate targeted studies are required to assess claim-to-action vulnerabilities in agentic NetOps workflows.
\end{itemize}

\section{Conclusion}
We executed a fully preregistered paired confirmatory trial comparing baseline free-text prompting with constrained-template prompting plus a single verifier-guided repair under shared\_initial\_candidate semantics. On a deterministic synthetic corpus with gpt-5-mini and deterministic\_netops\_validator\_v5, the guarded pipeline increased verifier pass rate from 149/150 to 150/150 (\ensuremath{\approx}0.67 pp), a numerically tiny and statistically non-significant difference (exact McNemar p = 1.0; 95\% bootstrap CI = [0.00, 0.02]). The principal scientific lesson is methodological: preregistered NetOps trials should include baseline calibration to avoid ceiling effects and preserve statistical power. Technically, our artifacts bound the marginal benefit of a one-repair guarded pipeline under shared\_initial\_candidate semantics; evaluating independent-sampling semantics, larger repair budgets, model diversity, and production-scale topologies remains important future work.

\section*{Disclosure Statement}
This study was executed autonomously after an autonomy lock under the archived master prompt (provenance/master\_prompt.txt, SHA-256 = 1872df1e\allowbreak{}1805d2d9\allowbreak{}6940456c\allowbreak{}a016bd66\allowbreak{}5d1d5196\allowbreak{}add77f5a\allowbreak{}cdf1582b\allowbreak{}b39b15ba). Preregistration identifier: study\_cand\_2026\_b2\_v1 (preregistration was frozen prior to observing outcomes and the preregistration record is archived). Declared model/tooling: declared model = gpt-5-mini (provider record: openai\_responses); execution adapter family = hosted\_netops\_gvr\_v1; deterministic validator = deterministic\_netops\_validator\_v5 (validator\_version = 5.0). Runtime sampling semantics (explicit): the archived model\_configuration.json records temperature = null (unset); because no numeric temperature or fixed random-seed was enforced, model outputs may be nondeterministic across independent API sessions. Autonomy and human role: a human Research Director selected the topic family and constrained the initial master prompt prior to the autonomy lock; after the lock the autonomous system performed literature review, study selection, preregistration, code generation, experiment execution, analysis, manuscript drafting, autonomous peer review, and revision. No human manually edited technical content, analyses, figures, captions, or references after the autonomy lock. Key execution facts reported in the paper (artifact-grounded): completed episodes = 300; complete paired tasks = 150; model\_calls\_used = 151; repair-stage invocations = 1. All runtime sampling metadata, provider-call traces, verifier inputs/verdicts, and analysis artifacts are archived in the execution bundle referenced by the execution manifest.

\begin{thebibliography}{99}
\bibitem{ref1} https://doi.org/10.1002/nem.2313
\bibitem{ref2} https://doi.org/10.1145/3786583.3786898
\bibitem{ref3} https://arxiv.org/abs/2604.22513
\bibitem{ref4} https://doi.org/10.1145/3605764.3623985
\bibitem{ref5} https://doi.org/10.1145/3656454
\bibitem{ref6} https://doi.org/10.1002/widm.70111
\end{thebibliography}

\end{document}
